\chapter{Grundlagen} \label{Grundlagen}
\section{Herleitung der regelungstechnischen Blockschaltbilder} \label{Section_Herleitung_Regelstrecke}
In diesem Abschnitt erfolgt die Herleitung der physikalischen Modelle für den \\ Tiefsetzsteller und den Hochsetzsteller.
Hierzu werden zunächst die jeweiligen\\  Systemgleichungen aufgestellt, anschließend normiert und mittels Laplace-Transformation in den Frequenzbereich überführt.
Vorhandene Nichtlinearitäten werden durch\\ geeignete Umrechnungsmechanismen modelliert und in die Systembeschreibung\\ integriert.
Das daraus abgeleitete Blockschaltbild bildet die Grundlage für die\\  anschließende Auslegung und Synthese des Regelungskonzepts.
\subsection{Tiefsetzsteller} \label{Section_lineare_Regelstrecke}
Zur Herleitung des regelungstechnischen Blockschaltbilds des Tiefsetzstellers dient das in Abbildung \ref{PhyM_TSS} dargestellte elektrische Ersatzschaltbild als Grundlage der weiteren Analyse.
Physikalische Größen werden mit "P" gekennzeichnet.
Nach der Normierung werden die Größen ohne "P" dargestellt.
\begin{figure}[h]
	\centering
	\includegraphics[width=8cm]{TSS_phy_s.pdf}
	\caption{Schaltbild eines Tiefsetzstellers}
	\label{PhyM_TSS}
\end{figure}\newpage \noindent
Die erste Systemgleichung des Systems ergibt sich aus der mittleren \\ Eingangsspannung $\overline{e}$, die durch das Tastverhältnis $\ILEAvar{\tau}{g}$ und die Eingangsspannung $\ILEAvar{U}{in,P}$ bestimmt wird.
\begin{align}
	\overline{e} &= \ILEAvar{\tau}{g} \cdot \ILEAvar{U}{in,P}
\end{align}
Die zweite Systemgleichung beschreibt den Zusammenhang zwischen der Spannung $\ILEAvar{u}{L,P}$ und dem Strom $\ILEAvar{i}{L,P}$ an der Spule $L$.
\begin{align}
	\ILEAvar{u}{L,P} &= L \cdot \derivative{\ILEAvar{i}{L,P}}{t}
\end{align} \noindent
Die dritte Systemgleichung ergibt sich aus dem Zusammenhang zwischen dem\\  Kondensatorstrom $\ILEAvar{i}{C,P}$ und der Kondensatorspannung $\ILEAvar{u}{C,P}$.
\begin{align}
	\ILEAvar{i}{C,P} &= C \cdot \derivative{\ILEAvar{u}{C,P}}{t}
\end{align} \noindent
Die vierte Systemgleichung leitet sich aus dem Ohmschen Gesetz ab. Dafür wird der Lastwiderstand $R$ betrachtet.
\begin{align}
	\ILEAvar{u}{C,P} &= R \cdot \ILEAvar{i}{R,P}
\end{align} \noindent
Um die fünfte Systemgleichung zu bestimmen, wird die Knotengleichung am\\  Ausgangsknoten betrachtet. Der Spulenstrom $\ILEAvar{i}{L,P}$ teilt sich in den Kondensatorstrom $\ILEAvar{i}{C,P}$ und den Laststrom $\ILEAvar{i}{R,P}$ auf.
\begin{align}
	\ILEAvar{i}{L,P} &= \ILEAvar{i}{C,P} + \underbrace{\ILEAvar{i}{R,P}}_{=z}
\end{align} \noindent
Die sechste Systemgleichung ergibt sich aus der Maschengleichung der\\  Eingangsspannung. Die Spulenspannung $\ILEAvar{u}{L,P}$ ist die Differenz zwischen der mittleren Eingangsspannung $\overline{e}$ und der Kondensatorspannung $\ILEAvar{u}{C,P}$.
\begin{align}
	\ILEAvar{u}{L,P} &= \overline{e} - \underbrace{\ILEAvar{u}{C,P}}_{=z}
\end{align} 
\newpage
\noindent
\textbf{Normierung der Systemgleichungen} \\
\noindent
Zur Vereinfachung der mathematischen Behandlung und zur Verallgemeinerung der Ergebnisse werden die physikalischen Größen auf $\ILEAvar{U}{N}$ und $\ILEAvar{I}{N}$ normiert:
\begin{align}
	\ILEAvar{T}{L} &= \frac{L \cdot \ILEAvar{I}{N}}{\ILEAvar{U}{N}} \\
	\ILEAvar{T}{C} &= \frac{C \cdot \ILEAvar{U}{N}}{\ILEAvar{I}{N}} \\
	\ILEAvar{u}{L} &= \frac{\ILEAvar{u}{L,P}}{\ILEAvar{U}{N}} \\
	\ILEAvar{u}{C} &= \frac{\ILEAvar{u}{C,P}}{\ILEAvar{U}{N}} \\
	\ILEAvar{u}{in} &= \frac{\ILEAvar{u}{in,P}}{\ILEAvar{U}{N}} \\
	e &= \frac{\overline{e}}{\ILEAvar{U}{N}} \\
	\ILEAvar{i}{L} &= \frac{\ILEAvar{i}{L,P}}{\ILEAvar{I}{N}} \\
	\ILEAvar{i}{C} &= \frac{\ILEAvar{i}{C,P}}{\ILEAvar{I}{N}} \\
	\ILEAvar{i}{R} &= \frac{\ILEAvar{i}{R,P}}{\ILEAvar{I}{N}} 
\end{align}
\newline \noindent
\textbf{Laplace-Transformation} \\
\noindent
Für das regelungstechnische Blockschaltbild werden die normierten \\ Systemgleichungen in den Laplace-Bereich transformiert:
\begin{align}
	\ILEAvar{u}{L} &= \ILEAvar{i}{L} \cdot p \ILEAvar{T}{L}  \\ 
	\ILEAvar{i}{C} &= \ILEAvar{u}{C} \cdot p \ILEAvar{T}{C}
\end{align} \newpage \noindent
\textbf{Regelungstechnisches Blockschaltbild} \\ \noindent
Das regelungstechnische Blockschaltbild des Tiefsetzstellers leitet sich direkt aus den normierten Systemgleichungen ab und ist in Abbildung \ref{RBS_TSS} dargestellt.
\begin{figure}[h]
	\centering
	\includegraphics[width=12cm]{TSS_RBS.pdf}
	\caption{Regelungstechnisches Blockschaltbild des Tiefsetzstellers}
	\label{RBS_TSS}
\end{figure} \\ \noindent
Das Blockschaltbild \ref{RBS_TSS} beschreibt das physikalische Modell des Tiefsetzstellers und wird für die Regelauslegung verwendet.


\subsection{Hochsetzsteller} \label{Section_Nichtlineare_Regelstrecke}
Wie in Abschnitt \ref{Section_lineare_Regelstrecke} leitet dieser Abschnitt das regelungstechnische Blockschaltbild des Hochsetzstellers anhand der normierten Systemgleichungen her.
Dazu betrachtet die Analyse das elektrische Ersatzschaltbild in Abbildung \ref{PhyM_HSS}. \\
\begin{figure}[h]
	\centering
	\includegraphics[width=8cm]{HSS_phyM.pdf}
	\caption{Schaltbild eines Hochsetzstellers}
	\label{PhyM_HSS}
\end{figure} \\ \noindent
Die erste Systemgleichung des Systems ergibt sich aus dem Zusammenhang zwischen der Spannung $\ILEAvar{u}{L,P}$  und dem Strom $\ILEAvar{i}{L,P}$ an der Spule $L$.
\begin{align}
	\ILEAvar{u}{L,P} &= L \cdot \derivative{\ILEAvar{i}{L,P}}{t}
\end{align} \noindent
Die zweite Systemgleichung beschreibt den Zusammenhang zwischen dem \\ Kondensatorstrom $\ILEAvar{i}{c,P}$ und der Ausgangsspannung $\ILEAvar{U}{out}$.
\begin{align}
	\ILEAvar{i}{C,P} &= C \cdot \derivative{\ILEAvar{U}{out}}{t}
\end{align}
Die dritte Systemgleichung ergibt sich aus dem Ohmschen Gesetz. Dafür wird der\\ Lastwiderstand $R$ betrachtet.
\begin{align}
	\ILEAvar{U}{out} &= R \cdot \ILEAvar{i}{R,P}
\end{align}
Um die vierte Systemgleichung zu bestimmen, wird der Diodenstrom $\ILEAvar{i}{D,p}$ betrachtet.
Dazu verwendet die Untersuchung das Steuergesetz des Hochsetzstellers und die\\  Beziehung zwischen Eingangs- und Ausgangsleistung.
\begin{align}
	\ILEAvar{U}{in} &= \left( 1 - \ILEAvar{\tau}{g} \right) \cdot \ILEAvar{U}{out} \\
	\ILEAvar{U}{in} \cdot\ILEAvar{i}{L,p} &= \ILEAvar{i}{D,p} \cdot \ILEAvar{U}{out} \\
	\ILEAvar{i}{D,P} &= \left( 1 - \ILEAvar{\tau}{g} \right) \cdot \ILEAvar{i}{L,p}  
\end{align}
\newpage
\noindent
Die letze Systemgleichungen wird am Ausgangsknoten des Hochsetzstellers\\  aufgestellt.
Der Diodenstrom $\ILEAvar{i}{D,P}$ teilt sich in den Kondensatorstrom $\ILEAvar{i}{C,P}$ und den\\ Lastwiderstandsstrom $\ILEAvar{i}{R,P}$ auf.
\begin{align}
	\ILEAvar{i}{D,P} &= \ILEAvar{i}{C,P} + \underbrace{\ILEAvar{i}{R,P}}_{=z} 
\end{align} \\ \noindent
\textbf{Normierung der Systemgleichungen} \\ \noindent
Zur Vereinfachung der mathematischen Behandlung und zur Verallgemeinerung der Ergebnisse werden die physikalischen Größen auf $\ILEAvar{U}{N}$ und $\ILEAvar{I}{N}$ normiert.
\begin{align}
	\ILEAvar{T}{L} &= \frac{L \ILEAvar{I}{N}}{\ILEAvar{U}{N}} \\
	\ILEAvar{T}{C} &= \frac{C \ILEAvar{U}{N}}{\ILEAvar{I}{N}} \\
	\ILEAvar{u}{L} &= \frac{\ILEAvar{u}{L,P}}{\ILEAvar{U}{N}} \\
	\ILEAvar{u}{C} &= \frac{\ILEAvar{u}{C,P}}{\ILEAvar{U}{N}} \\
	\ILEAvar{u}{in} &= \frac{\ILEAvar{u}{in,P}}{\ILEAvar{U}{N}} \\
	\ILEAvar{i}{L} &= \frac{\ILEAvar{i}{L,P}}{\ILEAvar{I}{N}} \\
	\ILEAvar{i}{C} &= \frac{\ILEAvar{i}{C,P}}{\ILEAvar{I}{N}} \\
	\ILEAvar{i}{R} &= \frac{\ILEAvar{i}{R,P}}{\ILEAvar{I}{N}} \\
	\ILEAvar{i}{D} &= \frac{\ILEAvar{i}{D,P}}{\ILEAvar{I}{N}} 
\end{align}
\textbf{Laplace-Transformation} \\ \noindent
Zur Bildung des regelungstechnischen Blockschaltbildes werden die normierten \\ Systemgleichungen in den Laplace-Bereich transformiert.
\begin{align}
	\ILEAvar{u}{L} &= \ILEAvar{i}{L} \cdot p \ILEAvar{T}{L}   \\
	\ILEAvar{i}{C} &= \ILEAvar{u}{C} \cdot p  \ILEAvar{T}{C}  
\end{align}
\newpage \noindent
\textbf{Regelungstechnisches Blockschaltbild} \\ \noindent
Somit ergibt sich das physikalische Modell des Hochsetzstellers in Abbildung \ref{PhyM_HSS} mit den Umrechenwerken (rote Umrandung) zu:
\begin{figure}[h]
	\centering
	\includegraphics[width=12cm]{HSS_RBS.pdf}
	\caption{Regelungstechnisches Blockschaltbild des Hochsetzstellers}
	\label{RBS_HSS}
\end{figure} \\ \noindent 
Ohne die Verwendung der Umrechenwerke hätte die Regelstrecke linearisiert werden müssen.
Dabei wäre die Regelung nur einsetzbar, wenn die Regelstrecke in der \\ Nähe des Arbeitspunktes betrieben wird.
In dieser Arbeit wird auf die Linearisierung des Hochsetzstellers verzichtet und stattdessen eine kaskadierte Regelung mit \\ Umrechenwerken verwendet.
Damit kann die Analyse direkt die Approximierfähigkeit des neuronalen Netzes im Vergleich einer linearen und nichtlinearen Regelstrecke \\ untersuchen. \cite{rothstielow2024leistungselektronik} \cite{rothstielow2024regelungstechnik}


